\section{An exact seed and event-driven hinge engine}
\label{sec:r15-hinge-core}
R15 selects an engine core with exact algebraic operands, finite expression
graphs, binary guard results and delayed state feedback. Geometry, branch
history and chart orientation are separately typed parts of its state. The
R14 electromagnetic, mechanical, thermal, clock and observation equations remain
the physical models reached through this interface. A geometric branch does
not itself supply force, energy, elapsed proper time or a detector reading.

\subsection{The golden-ratio pair is an exact arithmetic domain}
Let $\phi=(1+\sqrt5)/2$ and represent $a+b\phi$ by
$\operatorname{Phi}(a,b)$ with $a,b\in\Q$. Its representation is unique because
$\phi$ is irrational. The defining identity $\phi^2=\phi+1$ gives
\begin{align}
 (a,b)+(c,d)&=(a+c,b+d),\nonumber\\
 (a,b)(c,d)&=(ac+bd,ad+bc+bd),\nonumber\\
 \phi(a,b)&=(b,a+b),&
 \phi^{-1}(a,b)&=(b-a,a).
 \label{eq:r15-phi-pair}
\end{align}
The last two maps are exact inverses. For a nonzero pair its field inverse is
\begin{equation}
 (a+b\phi)^{-1}=\frac{(a+b)-b\phi}{a^2+ab-b^2}.
 \label{eq:r15-phi-inverse}
\end{equation}
The denominator is the product with the conjugate $a+b(1-\phi)$ and cannot
vanish for a nonzero rational pair. These are algebraic identities; a rational
approximation such as $1.618\ldots$ is not substituted for $\phi$.

For exact sign selection put $A=a+b/2$, $B=b/2$, so the value is
$A+B\sqrt5$. If $B=0$, use the sign of $A$. If $A=0$, use the sign of $B$.
If $A,B$ have the same sign, that is the answer. Otherwise
\begin{equation}
 \operatorname{sgn}(A+B\sqrt5)
 =\operatorname{sgn}(A)\operatorname{sgn}(A^2-5B^2).
 \label{eq:r15-phi-sign}
\end{equation}
The last comparison is rational and is justified by comparing two positive
magnitudes before squaring. Thus this domain admits exact comparisons without
evaluating a rounded square root. Transcendental expressions outside the
selected domain retain the inherited unresolved-result convention.

\subsection{One-bit planes retain arithmetic carries and the whole expression}
Each rational coefficient has an integer numerator and a positive denominator;
canonical reduction uses their greatest common divisor. Integer magnitudes,
signs, lengths, denominators, operator tags and references all belong to the
source. For a group of $w$ serialized operands, their bit plane $k$ can be
packed as
\begin{equation}
 B_k=\sum_{j=0}^{w-1}b_{j,k}2^j,\qquad
 b_{j,k}=(B_k\mathbin{\gg}j)\mathbin{\&}1.
 \label{eq:r15-bitplanes}
\end{equation}
Declared lengths and padding make this a bijection on the finite payload.
One-bit encoding changes the layout, not the number of mathematical degrees
of freedom. A variable-length exact coefficient is not reduced to one scalar
bit or one fixed-width word.

For nonnegative binary integers the full adder is explicitly
\begin{align}
 c_0&=0,&s_i&=a_i\oplus b_i\oplus c_i,\nonumber\\
 c_{i+1}&=(a_i\land b_i)\lor[c_i\land(a_i\oplus b_i)].
 \label{eq:r15-carry}
\end{align}
The final carry is retained. Subtraction uses an explicit signed-integer
operation or sufficient sign extension and two's-complement addition, with
the complete signed result retained. Fixed logical words still use their
declared modulo-$2^w$ arithmetic, a different type. The pair update
$(a,b)\mapsto(b,a+b)$ therefore needs addition with carry; a shift or XOR alone
does not implement multiplication by $\phi$. Numerator/denominator arithmetic
likewise retains all carries and any required additional limbs.

The seed's finite directed acyclic graph stores literal and input nodes,
selected arithmetic, affine points and guard bodies. A literal is an exact
operand; an input node binds one immutable sample. Operand order and grouping
are part of the graph, including sequential arithmetic folds. Lazy branches
evaluate only their demanded operands. Shared subexpressions may be cached
under the same operand identities without changing the denotation.

The source's literal ``zero-order'' convention is left-to-right sequencing
with uniform operator precedence. For the explicit sequence
$(v_0,(\circ_1,v_1),\ldots,(\circ_m,v_m))$, its body is
\begin{equation}
 a_0=v_0,\qquad a_i=a_{i-1}\circ_i v_i\quad(1\le i\le m),
 \qquad \operatorname{Fold}=a_m.
 \label{eq:r15-sequential-fold}
\end{equation}
Every intermediate value is type checked. Explicitly grouped subexpressions
are operand nodes with their own bodies. Thus an ungrouped addition followed
by multiplication is folded in that order; a compiler must not impose the
usual multiplication-before-addition parse on this source convention. The
stored AST records the selected grouping unambiguously.

``Zero-order'' here says nothing about derivative order, absence of dynamics
or $O(1)$ computational complexity. Work depends on visited graph nodes,
operand bit lengths, comparisons and generated descendants. An exact
expression may still be expensive to evaluate. Independently selected physical
evolution retains its actual body; the fold does not replace the inherited
\code{ROOT}, \code{ODE} or \code{DIFF} domains when those are requested.

\subsection{Guard values, boundary ownership and an accepted event}
A geometric guard is a signed defining function $g$ whose selected inside is
$g<0$. A plane may use $g(x)=n\cdot x-c$, with $n\ne0$; a sphere may use
\begin{equation}
 g_{\rm sphere}(x)=\|x-a\|^2-R^2,\qquad R>0.
 \label{eq:r15-sphere-guard}
\end{equation}
The latter is a squared polynomial guard, not a distance in metres.
It has the sphere's exact sign while avoiding a square root in that predicate.
The sign theorem for rational $\phi$ pairs applies when all active operands
belong to that domain.

Let $\kappa\in\B$ record the current local coorientation and set
$g_c=(-1)^\kappa g$. The guard bit is $b_g=[g_c<0]$. Equality is a separate
\code{BOUNDARY} result by default, with no transition commit. An application
may explicitly assign boundary ownership to inside or outside. An unavailable
demanded sample is \code{MISSING\_INPUT}; an unsupported demanded operator is
\code{UNKNOWN}. Neither status is a false guard.

The event record binds an immutable identity and content to the selected
sample, seam and event-order inputs. Physical acquisition additionally
supplies its declared timestamp and availability binding; a logical sequence
counter alone is not elapsed physical time. A new accepted event supplies a hinge pulse $h\in\B$:
it can be a side change after a defined nonboundary baseline, an explicitly
declared crossing, or a declared nontrivial seam transition. An initial sample
can establish the baseline without a pulse. A repeated identical identity is
not a new pulse; reusing an identity with changed content is invalid. Continuous
trajectory crossings not represented by this acquisition/event policy are not
proved detected merely because sampled endpoints have the same sign.

\subsection{Every hinge crossing and a reversing seam are different bits}
Let $p$ be hinge parity, $\kappa$ coorientation parity, and
$\epsilon\in\B$ identify whether the admitted seam reverses the selected
local sign. The next parities are
\begin{equation}
 p^+=p\oplus h,\qquad
 \kappa^+=\kappa\oplus(h\land\epsilon).
 \label{eq:r15-two-parities}
\end{equation}
Thus every accepted hinge pulse changes $p$, whereas an orientation-preserving
hinge leaves $\kappa$ unchanged. The determinant $-1$ of the coefficient-pair
map in equation~\eqref{eq:r15-phi-pair} concerns an algebraic basis; it is not
by itself a physical chart-orientation theorem.

The inherited two independently stored ASA/NA mask sets act in their declared
order and produce the word $y$ from the common old snapshot. Define
$b_{\rm lat}=\operatorname{bit}(y,d)$ at a reserved drive lane $d$.
It is this post-mask bit, rather than an unconditionally substituted raw
guard, that the selected latch records. The selected default drive adapter
sets
\begin{equation}
 J_d=h\land b_{\rm lat},\qquad K_d=h\land\neg b_{\rm lat}.
 \label{eq:r15-drive-lane}
\end{equation}
On a distinct optional seam lane $s$ set $J_s=K_s=h\land\epsilon$.
Every other lane retains its base $J,K$ body. Assemble the word masks and use
the same synchronous JK identity on all lanes:
\begin{equation}
 q^+=(J^*\land\neg_w q)\lor(\neg_w K^*\land q).
 \label{eq:r15-jk}
\end{equation}
Equation~\eqref{eq:r15-drive-lane} holds when $h=0$ and sets the drive lane to
$b_{\rm lat}$ when $h=1$; the seam lane toggles exactly on a reversing pulse.
Custom word equations must retain this body to claim this particular drive
adapter; an arbitrary user-supplied $J,K$ expression need not implement it.
Disjoint allocation prevents either adapter from overwriting the other's
semantics. This is local lane preservation for one step: later base expressions
may read those changed bits, so it is not a global noninterference claim.

Guard, mask, parity, chart-lift, consumed-event and word outputs are computed
from one old state and accepted record, then committed together. A seam's new
raw guard must be paired with the matching new coorientation before comparing
physical sides; applying old $\kappa$ to a new reversed guard would invent a
crossing. No partial state is committed for an unresolved demanded result.
Physical time and an external plant can continue while a digital event waits.

\subsection{Seam compatibility, winding and a graph beyond the branch tree}
For overlapping charts $U_i,U_j$, declare an invertible transition $T_{ij}$
and a sign $\sigma_{ij}=(-1)^{\epsilon_{ij}}$. The selected exact local guard
compatibility is
\begin{equation}
 g_j(T_{ij}x)=\sigma_{ij}g_i(x).
 \label{eq:r15-sign-compatibility}
\end{equation}
For general defining functions a positive factor could preserve sign, but it
must then be an explicit different contract; exact SDF transport has its
stronger metric conditions in Section~\ref{sec:r15-spatial-calculus}.
An input record asserting this scalar equality does not by itself prove that
both arguments denote the same physical point. That premise is supplied by
the actual chart map and coordinate binding.

On a triple overlap the maps and signs obey
\begin{equation}
 T_{ik}=T_{jk}\circ T_{ij},\qquad
 \sigma_{ik}=\sigma_{jk}\sigma_{ij}.
 \label{eq:r15-cocycle}
\end{equation}
When using lifted periodic coordinates these equalities are interpreted with
the explicitly retained deck/winding action, not by dropping a full turn.
For example, $\vartheta=\bar\vartheta+2\pi w$ with $w\in\Z$ and
$\bar\vartheta\in[0,2\pi)$ retains the complete angle. Normalizing its address
updates $w$; the residue alone is insufficient to invert the lift. The inherited
half-turn and reflective-Klein chart actions remain separately named maps.

For a closed seam walk $C$, the product $\prod_{(i,j)\in C}\sigma_{ij}$ is
its sign holonomy. If it is $-1$, a globally single-valued choice of local
signs cannot trivialize that loop. This is a statement about the specified
chart/sign bundle, not a proof that an arbitrary embedded solid has a globally
defined signed distance or a particular physical Klein-bottle geometry.
A branch tree alone has no cycle; repeated bit flips along a tree cannot
establish such holonomy. The engine therefore distinguishes its expansion
tree from an additional seam graph containing declared return and cross edges.
Every return is an explicit state transition at a later accepted event. It
does not impose a self-consistent eigenvector or an instantaneous circular AST.

\subsection{CUDA realization: immutable source and derived VRAM epochs}
The native realization has two kinds of device data. The authoritative source
is the immutable packed operator/operand payload and its version identity.
Expansion products, bounds, classifications and work queues are derived VRAM
data. Their key includes source identity, state/event epoch, parameters and
the represented interval. Rebuilding a cache is permitted only if the same
declared semantic result or a certified enclosure is preserved.

For the raw-word path, a linear texture object exposes unsigned integer
elements and uses \code{tex1Dfetch} with integer indices and
\code{cudaReadModeElementType}. No filtered interpolation, normalization,
color conversion or texture-coordinate wrapping is a semantic operation on
the source bits. Bounds, alignment and resource size are checked explicitly.
The relevant descriptor semantics are documented in
\href{https://docs.nvidia.com/cuda/cuda-runtime-api/cuda_runtime_api/group__CUDART__TEXTURE__OBJECT.html}
{HC1: NVIDIA CUDA Runtime API, Texture Object Management}.

A producer may expand the read-only seed into a separate output allocation.
A completed kernel and explicit stream/event dependency precede a consumer
reading that output as a texture in the next epoch. No kernel reads through
the texture path from storage that it concurrently modifies through global
writes. The reason is concrete: NVIDIA specifies that the texture cache is
not coherent with same-launch global writes; see
\href{https://docs.nvidia.com/cuda/cuda-c-best-practices-guide/#texture-memory}
{HC2: CUDA C++ Best Practices Guide, Texture Memory}.
Double buffering or another explicit ownership schedule can satisfy this
contract. A device pointer or a launch-order assumption across unsynchronized
streams is not a completed dependency.

The device lowering declares its supported subset. A kernel performing native
bounded-width integer word operations is not thereby a complete arbitrary-
precision rational-$\phi$ evaluator. A kernel consuming projected IEEE
coordinates is not thereby the exact AST evaluator. Each supported predicate
either executes its exact body, or returns a proved enclosure/decision with a
defined exact fallback for unresolved cases. Unsupported bodies preserve the
source and route to a supported evaluator; they are not silently rounded into
different semantics. Correctness checks establish the tested subset, and
benchmarks separately report its time, transfers and resource use.

Hardware chooses texture-cache residency, replacement and scheduling. The
formal contract guarantees neither a cache hit nor a speed advantage over
ordinary global loads. Seed expansion in VRAM is a realizable dataflow and
implementation choice whose capacity and throughput are measured for the
actual GPU. It does not mean an infinite tree resides in finite memory or that
one-bit storage alone computes its operators. Self-reference remains the
explicit delayed event/state recurrence above.
